\section{Введение}
\label{sec:Chapter0} \index{Chapter0}

% -----------------------------------------------------------------------
% 1. АКТУАЛЬНОСТЬ И КОНТЕКСТ
% -----------------------------------------------------------------------

Современные промышленные рекомендательные системы работают в условиях
непрерывной обратной связи: модель формирует выдачу, пользователь реагирует
на неё, журналы взаимодействий поступают в обучающую выборку следующей
итерации дообучения. Такой режим эксплуатации принципиально отличается от
стандартного предположения машинного обучения о независимой и одинаково
распределённой выборке (i.i.d.): распределение обучающих данных на каждом
шаге зависит от параметров модели на предыдущем шаге. Это явление принято
называть \textit{скрытой петлёй обратной связи} (hidden feedback loop).

Проблема hidden feedback loops в системах машинного обучения изучается
как в контексте общей теории repeated learning, так и применительно
к конкретным классам задач: ранжирование, фильтрация контента,
персонализация новостных лент. В рекомендательных системах эта петля
замкнута особенно тесно: пользователь видит только то, что система
рекомендует, а система обучается только на том, что пользователь увидел и
на что отреагировал. Следовательно, области пространства объектов,
недоступные пользователю в момент времени $t$, выпадают из обучающей
выборки момента $t{+}1$. Накопленный эффект многократного повторения этого
цикла может привести к систематическому сужению выдачи и, в пределе, к
формированию так называемых \textit{эхо-камер}~--- режима, при котором
разные кластеры пользователей получают диверсифицированные, но внутренне
однородные рекомендации без пересечения с контентом других кластеров.

Геометрическая интерпретация этого явления связана с изменением структуры
пользовательских предпочтений в процессе repeated learning. Если
рассматривать состояние системы как распределение вероятностей в
пространстве пользователей, то каждая итерация переобучения действует как
оператор, трансформирующий это распределение. При определённых условиях
такой оператор оказывается сжимающим: внутрикластерное разнообразие
убывает, а межкластерные границы становятся резче. Анализ условий, при
которых это сжатие неизбежно, и является центральным предметом настоящей
работы.

% -----------------------------------------------------------------------
% 2. ОБЪЕКТ И ПРЕДМЕТ ИССЛЕДОВАНИЯ
% -----------------------------------------------------------------------

\textbf{Объектом исследования} является рекомендательная система,
функционирующая в режиме онлайн-переобучения с обратной связью по
пользовательским реакциям.

\textbf{Предметом исследования} служит геометрия распределения
рекомендаций в пространстве пользователей в условиях repeated learning:
в частности, условия затухания внутрикластерной дисперсии выдачи и
возникновения эхо-камерного режима.

% -----------------------------------------------------------------------
% 3. ЦЕЛЬ И ЗАДАЧИ
% -----------------------------------------------------------------------

\textbf{Цель работы}~--- построить теоретическую модель, описывающую
долгосрочную динамику персонализации в рекомендательной системе с обратной
связью, и установить условия, при которых эта динамика приводит к
деградации микроперсонализации.

Для достижения поставленной цели необходимо решить следующие задачи:
\begin{enumerate}
    \item Формализовать рекомендательную систему с повторяющимся
          переобучением как дискретную динамическую систему на пространстве
          распределений пользователей и объектов.
    \item Ввести кластерную структуру аудитории и получить локальную
          линеаризацию оператора переобучения вблизи стационарного
          распределения.
    \item Исследовать спектральные свойства эффективного оператора
          переобучения и установить условия его сжимающего действия на
          внутрикластерные ковариации.
    \item Доказать теорему о коллапсе персонализации: при выполнении условий
          сжатия распределение рекомендаций сходится к $K$-модальной
          комбинации дельта-мер.
    \item Верифицировать теоретические выводы численно на синтетических и
          реальных данных, измерив динамику внутрикластерного разнообразия
          при различных режимах переобучения.
    \item Проанализировать влияние управляющих параметров~--- уровня
          exploration и силы регуляризации~--- на предельный режим системы
          и установить достаточные условия сохранения ненулевого разнообразия.
\end{enumerate}

% -----------------------------------------------------------------------
% 4. МЕТОДЫ ИССЛЕДОВАНИЯ
% -----------------------------------------------------------------------

В работе используются методы теории динамических систем на пространствах
мер (операторный подход, теория сжимающих отображений), линейной алгебры
и спектральной теории (анализ спектрального радиуса, сингулярное
разложение), теории вероятностей (ковариационная динамика, слабая
сходимость мер), а также вычислительные методы: моделирование repeated
recommendation loop на синтетических данных и датасете MovieLens.

% -----------------------------------------------------------------------
% 5. НАУЧНАЯ НОВИЗНА И ПРАКТИЧЕСКАЯ ЗНАЧИМОСТЬ
% -----------------------------------------------------------------------

\textbf{Научная новизна} работы состоит в следующем. Во-первых, предложена
кластерная динамическая модель рекомендательной системы, явно учитывающая
зависимость обучающих данных от политики рекомендаций. Во-вторых, получены
спектральные условия, при которых оператор переобучения является сжимающим
по внутрикластерным ковариациям. В-третьих, доказана теорема о сходимости
распределения рекомендаций к $K$-модальной эхо-камере и установлена
зависимость скорости этой сходимости от параметров модели.

\textbf{Практическая значимость} определяется тем, что полученные
теоретические условия допускают непосредственную проверку по наблюдаемым
статистикам системы и могут служить основой для разработки метрик
мониторинга риска коллапса персонализации в промышленных рекомендательных
системах.

% -----------------------------------------------------------------------
% 6. СТРУКТУРА РАБОТЫ
% -----------------------------------------------------------------------

Работа организована следующим образом. В~главе~\ref{sec:Chapter1}
рекомендательная система формализуется как repeated learning process и
ставится основной исследовательский вопрос. Глава~\ref{sec:Chapter2}
посвящена построению математической модели: вводятся кластерная структура
аудитории, локальная линеаризация рекомендательной политики и ковариационная
динамика выдачи. В~главе~\ref{sec:Chapter3} формулируются и доказываются
основные теоремы~--- о спектральном сжатии и коллапсе персонализации.
Глава~\ref{sec:Chapter4} описывает вычислительный эксперимент и анализирует
соответствие численных результатов теоретическим предсказаниям.
Глава~\ref{sec:Chapter5} рассматривает управляемую версию модели и анализирует
условия предотвращения коллапса через exploration и регуляризацию. В
заключении (глава~\ref{sec:Chapter6}) формулируются итоги работы и
направления дальнейших исследований.
